\section{Introduction}
\label{sec:introduction}

Multi-protocol DeFi lending allocation has been studied either as a
rate-forecasting problem on hourly-aggregated rate series
\cite{baude2025optimal} or as a control-theoretic problem with
reinforcement learning on the protocol-design side
\cite{bertucci2025rl, xu2023autogov, agilerate2024}.  Both
formulations share an implicit assumption: that hourly resolution is
sufficient for an allocator to extract economically meaningful
edge.  We argue that this assumption is empirically wrong.
Lending-rate crossovers between Ethereum L1 protocols happen at the
speed of pool deposits and withdrawals --- per-block ($\sim 12$\,s
post-PoS), not per-hour --- because the rate is a deterministic
function of realised utilization, and utilization updates exactly
once per on-chain interaction \cite{werner2022sok}.
Aggregating to hourly bars destroys the inter-block microstructure
where the allocator's competitive edge actually lives.

\paragraph{Thesis.}
We argue that DeFi lending allocation is structurally analogous to
high-frequency trading microstructure as described by
\cite{mackenzie2021}, and that the four-class HFT signal taxonomy of
that work transposes onto the multi-protocol lending problem with a
clean mapping.  Concretely: the cross-protocol rate spread is itself
the decision variable (the dominant signal); mempool-visible
transactions are the analog of ``order-book dynamics''; lead rates
from Maker DSR and Curve 3pool stand in for the ``futures lead''
class; and the on-chain gas regime plus stablecoin peg deviations
play the role of ``related instruments''.  Under this lens the
natural decision frame is event-time, gas-aware, with three nested
policy tiers --- gas-aware threshold (T1), Ornstein--Uhlenbeck
optimal stopping with switching cost (T2), and Cox proportional-
hazards regression on the signal vector (T3) --- each grounded in a
separate microstructure or quantitative-finance source book
\cite{ohara1995, krause2005, kissell2014, lopezdeprado2018,
mackenzie2021}.  This formulation extends the published 4-factor
MCDM baseline of \cite{solovev2026vault} from hourly cliff-edge
switching to per-block gas-aware switching, and from a 2-way (Aave
$+$ Compound) panel to the 6-protocol matrix that covers
$\sim 67\%$ of total Ethereum-L1 USDC-lending TVL.

\paragraph{Empirical finding.}
On a $\NBlocksFull$ per-block panel covering
\DataStart{}--\DataEnd{} across $\NProtocols$ Ethereum USDC lending
venues (\ProtocolList) plus three passive-hold benchmarks
(Compound V3, Spark, Fluid), T1 beats passive Aave V3 hold under
paired monthly Sharpe bootstrap
($\Delta\textsc{Sharpe} = \HOneAuxDelta$, 95\,\% CI
$[\HOneAuxCIlow, \HOneAuxCIhigh]$, $p = \HOneAuxPvalue$); the
sophisticated-retrain T3 (Cox PH on the F1$+$F3 design matrix)
closes pre-registered H1c with $\Delta\textsc{APY} = +7.03$\,bp
over T1 ($p = 0.015$) on the $N\!=\!6$ walk-forward paired
bootstrap.  Two orders of magnitude more rebalance opportunities
than the hourly baseline of \cite{solovev2026vault} expresses.

\paragraph{Contributions.}
\begin{itemize}
  \item \textbf{An event-time, gas-aware multi-protocol DeFi
    lending allocator with three nested decision-policy tiers
    (T1/T2/T3), each grounded in a distinct microstructure or
    quantitative-finance source.}  We give the closed-form gas-cost
    crossover inequality that locks T1's threshold, the
    OU--Bellman analytical boundary that closes T2, and the Cox
    proportional-hazards specification with the triple-barrier
    label of \cite{lopezdeprado2018} for T3.

  \item \textbf{The four-class HFT signal taxonomy of
    \cite{mackenzie2021} (Table 3.2) operationalised on DeFi
    lending.}  F1 (lead) $\leftarrow$ Maker DSR rate plus lagged
    spreads to the top-protocol APR; F3 (fragmentation)
    $\leftarrow$ the cross-protocol rate spread itself, the
    dominant signal and the F3-only T3 fit's sole driver;
    F4 (related) $\leftarrow$ ETH/USD, gas regime, USDC peg
    deviations; F2 (order-book dynamics) deferred per the
    explicit risk register pending paid mempool coverage.

  \item \textbf{L\'opez de Prado AFML methodology applied to DeFi
    lending}: triple-barrier labels (24-h horizon), purged $k$-fold
    CV with embargo $\approx 5.4$ days per
    \cite[Ch.\,7]{lopezdeprado2018}, sample-uniqueness weighting,
    Deflated Sharpe Ratio gate
    \cite[Ch.\,14.7.3]{lopezdeprado2018}.

  \item \textbf{The first per-block-resolution head-to-head matrix
    of seven allocation policies on a multi-month Ethereum mainnet
    panel} (four baselines plus T1/T2/T3), with walk-forward
    $N \!\times\! M$ paired bootstrap as the binding inference and
    a regime-conditional Q1/Q2 breakdown that exposes which tier
    dominates which volatility regime.

  \item \textbf{A production-grade reference agent + full
    reproducibility envelope.}  The same \texttt{decision/} Python
    package is consumed \emph{bit-identically} by an off-chain
    agent (per-block \texttt{newHeads} subscription, Flashbots
    private mempool --- the asymmetric speed-bump of
    \cite[pp.\,200--203]{mackenzie2021};
    \cite{heimbach2023pbs} survey the
    defence space).  The 6-way panel ($\sim$3.9M rows), Cox
    sidecars, walk-forward equity parquets, and Institutional
    Dossier are public on HuggingFace + GitHub.
\end{itemize}

\paragraph{Outline.}
\Cref{sec:background} reviews DeFi lending mechanics, the HFT
canon we borrow from, and recent DeFi-microstructure / MEV work
\cite{capponi2025mev, lehar2024fragility}.
\Cref{sec:methodology} specifies the three-tier ladder.
\Cref{sec:empirical} reports the seven-policy matrix and the
binding walk-forward bootstrap.  \Cref{sec:discussion} positions
the result against the HFT canon and the cross-domain hinge;
\Cref{sec:limitations,sec:conclusion} state scope and next steps.
